% ─────────────────────────────────────────────────────────────────────────────
% LAMPIRAN
% Cantumkan berkas-berkas lampiran yang relevan terhadap proposal (misalnya:
% surat izin penelitian, instrumen/kuesioner, contoh data, kode program).
% Penomoran halaman mengikuti angka romawi bagian awal (diatur di file utama).
% Format judul: kiri atas, nomor lampiran cetak tebal, dan dicantumkan di
% Daftar Isi dengan pola:
%   \addcontentsline{toc}{section}{Lampiran N Judul Lampiran}
% Hapus contoh lampiran di bawah dan ganti dengan lampiran Anda.
% ─────────────────────────────────────────────────────────────────────────────
% ─────────────────────────────────────────────────────────────────────────────
\newpage
\phantomsection
\addcontentsline{toc}{chapter}{LAMPIRAN}
\noindent\textbf{Lampiran 1} Hasil Pemeriksaan \textit{AI Detector}%
\phantomsection
\addcontentsline{toc}{section}{Lampiran 1 Hasil Pemeriksaan AI Detector}

\vspace{1em}
Sisipkan bukti hasil pemeriksaan naskah menggunakan perangkat deteksi AI (\textit{AI detector}) premium yang diinformasikan oleh Koordinator Tugas Akhir, misalnya tangkapan layar atau laporan PDF. Tingkat indikasi penggunaan AI tidak boleh melebihi 20\%; apabila melebihi, naskah wajib direvisi sebelum diajukan untuk proses penilaian.
% Contoh: \includegraphics[width=\linewidth]{Gambar/hasil-ai-detector.png}

% ─────────────────────────────────────────────────────────────────────────────
\newpage
\noindent\textbf{Lampiran 2} Contoh Lampiran Kode Program%
\phantomsection
\addcontentsline{toc}{section}{Lampiran 2 Contoh Lampiran Kode Program}

\vspace{1em}
Kode program yang panjang sebaiknya diletakkan di lampiran. Gunakan perintah \verb|\lstinputlisting| untuk menyisipkan langsung dari file, misalnya:
\par\noindent\verb|\lstinputlisting[language=Python]{Kode/nama_file.py}|

\noindent Untuk potongan kode di dalam bab, gunakan lingkungan \texttt{lstlisting} dengan opsi \texttt{caption} agar tercantum dalam Daftar Source Code (lihat contoh pada Bab~III).

% ─────────────────────────────────────────────────────────────────────────────
\newpage
\noindent\textbf{Lampiran 3} Contoh Cara Menyusun Banyak Gambar%
\phantomsection
\addcontentsline{toc}{section}{Lampiran 3 Contoh Cara Menyusun Banyak Gambar}

\vspace{1em}
Gunakan lingkungan \texttt{subfigure} dari paket \texttt{subcaption} untuk menempatkan beberapa gambar berdampingan dalam satu \texttt{figure}. Pastikan total lebar subfigure tidak melebihi \verb|\linewidth|. Contoh struktur kode untuk dua gambar berdampingan:

\begin{lstlisting}[language={[LaTeX]TeX}]
\begin{figure}[H]
  \centering
  \begin{subfigure}{.45\linewidth}
    \centering
    \includegraphics[width=\linewidth]{Gambar/gambar1.png}
    \caption{Keterangan Gambar 1}
  \end{subfigure}
  \hfill
  \begin{subfigure}{.45\linewidth}
    \centering
    \includegraphics[width=\linewidth]{Gambar/gambar2.png}
    \caption{Keterangan Gambar 2}
  \end{subfigure}
  \caption{Judul Keseluruhan Gambar}
  \label{fig:multi}
\end{figure}
\end{lstlisting}
